% GENERATED FILE. Edit canonical lessons, then run npm run generate:books.
% canonical-source-hash: fnv1a64:5143feeb1ca3aced
% canonical-lessons: MR-C311-vartahar, MR-C311-sampadak, MR-C311-lekh, MR-C311-agralekh, MR-C311-prakasit-karne, MR-R311-review-304-311

\chapter{Who Makes the News}
\label{ch:mr-who-makes-the-news}
\hlchaptermodality{311}

\begin{chapteropening}
\textbf{By the end of this chapter:} \emph{I can say who reports and edits the news.}

Complete the last lesson of chapter 311: I can say who reports and edits the news.
\end{chapteropening}

\section[\texorpdfstring{\mr{वार्ताहर} (vārtāhar)}{वार्ताहर (vārtāhar)}]{\mr{वार्ताहर} (vārtāhar) — a reporter; a correspondent}

\label{lesson:MR-C311-vartahar}



\begin{quote}
Before the new one: say the Marathi for a headline, then the Marathi for to print.
\end{quote}

\subsection*{You'll want to know: \mr{वार्ताहर}}
\textbf{\mr{वार्ताहर}} — \emph{vārtāhar} — \textquotedblleft{}a reporter; a correspondent\textquotedblright{}.

\begin{quote}
\mr{वार्ताहराने} \mr{अपघाताची} \mr{माहिती} \mr{फोनवर} \mr{दिली}.
\end{quote}

\begin{quote}
\emph{vārtāharāne apghātācī māhitī phonvar dilī.}
\end{quote}

\begin{quote}
The reporter gave the details of the accident over the phone.
\end{quote}

Masculine.

Built from \mr{वार्ता}, news, and -\mr{हर}, carrier.

\begin{cousinweb}
\textbf{\mr{वार्ताहर}} (\emph{vārtāhar}) is a tatsama, taken from Sanskrit \emph{vārttāhara}, 'a carrier of tidings, a messenger' (Monier-Williams).
\end{cousinweb}

\begin{grammarlens}[title={What you've built}]
One more word for telling what happened.
\end{grammarlens}

\subsection*{Your turn}
\begin{itemize}
\raggedright
  \item \emph{Say it:} \emph{vārtāhar}
  \item \emph{Say it:} \emph{vārtāhar}, once more
  \item \emph{Say it:} \emph{vārtāharāne apghātācī māhitī phonvar dilī.} — The reporter gave the details of the accident over the phone.
  \item \emph{Recall:} say the Marathi for a headline, then the Marathi for to print, then say \emph{vārtāhar} again
\end{itemize}

\begin{tcolorbox}[breakable,colback=teal!4,colframe=teal!35!black,title={Before you move on}]
What does \mr{वार्ताहर} mean? (\textquotedblleft{}A reporter; a correspondent\textquotedblright{}.) Say it once more.
\end{tcolorbox}
\section[\texorpdfstring{\mr{संपादक} (sampādak)}{संपादक (sampādak)}]{\mr{संपादक} (sampādak) — an editor}

\label{lesson:MR-C311-sampadak}



\begin{quote}
Before the new one: say the Marathi for to print, then the Marathi for a reporter.
\end{quote}

\subsection*{You'll want to know: \mr{संपादक}}
\textbf{\mr{संपादक}} — \emph{sampādak} — \textquotedblleft{}an editor\textquotedblright{}.

\begin{quote}
\mr{संपादकांनी} \mr{लेखात} \mr{काही} \mr{बदल} \mr{केले}.
\end{quote}

\begin{quote}
\emph{sampādakānnī lekhāt kāhī badal kele.}
\end{quote}

\begin{quote}
The editor made some changes to the article.
\end{quote}

Masculine.

The respectful plural \mr{संपादकांनी}, by the editor.

\begin{cousinweb}
\textbf{\mr{संपादक}} (\emph{sampādak}) is a tatsama, taken from Sanskrit \emph{sampādaka}, 'one who brings about, accomplishes' (Monier-Williams); the sense 'editor' is modern.
\end{cousinweb}

\begin{grammarlens}[title={What you've built}]
One more word for telling what happened.
\end{grammarlens}

\subsection*{Your turn}
\begin{itemize}
\raggedright
  \item \emph{Say it:} \emph{sampādak}
  \item \emph{Say it:} \emph{sampādak}, once more
  \item \emph{Say it:} \emph{sampādakānnī lekhāt kāhī badal kele.} — The editor made some changes to the article.
  \item \emph{Recall:} say the Marathi for to print, then the Marathi for a reporter, then say \emph{sampādak} again
\end{itemize}

\begin{tcolorbox}[breakable,colback=teal!4,colframe=teal!35!black,title={Before you move on}]
What does \mr{संपादक} mean? (\textquotedblleft{}An editor\textquotedblright{}.) Say it once more.
\end{tcolorbox}
\section[\texorpdfstring{\mr{लेख} (lekh)}{लेख (lekh)}]{\mr{लेख} (lekh) — an article}

\label{lesson:MR-C311-lekh}



\begin{quote}
Before the new one: say the Marathi for a reporter, then the Marathi for an editor.
\end{quote}

\subsection*{You'll want to know: \mr{लेख}}
\textbf{\mr{लेख}} — \emph{lekh} — \textquotedblleft{}an article\textquotedblright{}.

\begin{quote}
\mr{तिचा} \mr{लेख} \mr{आजच्या} \mr{वृत्तपत्रात} \mr{आला}.
\end{quote}

\begin{quote}
\emph{ticā lekh ājacyā vṛttapatrāt ālā.}
\end{quote}

\begin{quote}
Her article appeared in today's newspaper.
\end{quote}

Masculine.

\mr{लेख} \mr{लिहिणे}, to write an article.

\begin{cousinweb}
\textbf{\mr{लेख}} (\emph{lekh}) is a tatsama, taken from Sanskrit \emph{lekha}, 'a line, a writing' (Monier-Williams).
\end{cousinweb}

\begin{grammarlens}[title={What you've built}]
One more word for telling what happened.
\end{grammarlens}

\subsection*{Your turn}
\begin{itemize}
\raggedright
  \item \emph{Say it:} \emph{lekh}
  \item \emph{Say it:} \emph{lekh}, once more
  \item \emph{Say it:} \emph{ticā lekh ājacyā vṛttapatrāt ālā.} — Her article appeared in today's newspaper.
  \item \emph{Recall:} say the Marathi for a reporter, then the Marathi for an editor, then say \emph{lekh} again
\end{itemize}

\begin{tcolorbox}[breakable,colback=teal!4,colframe=teal!35!black,title={Before you move on}]
What does \mr{लेख} mean? (\textquotedblleft{}An article\textquotedblright{}.) Say it once more.
\end{tcolorbox}
\section[\texorpdfstring{\mr{अग्रलेख} (agralekh)}{अग्रलेख (agralekh)}]{\mr{अग्रलेख} (agralekh) — an editorial}

\label{lesson:MR-C311-agralekh}



\begin{quote}
Before the new one: say the Marathi for an editor, then the Marathi for an article.
\end{quote}

\subsection*{You'll want to know: \mr{अग्रलेख}}
\textbf{\mr{अग्रलेख}} — \emph{agralekh} — \textquotedblleft{}an editorial\textquotedblright{}.

\begin{quote}
\mr{आजचा} \mr{अग्रलेख} \mr{पाण्याच्या} \mr{प्रश्नावर} \mr{आहे}.
\end{quote}

\begin{quote}
\emph{ājcā agralekh pāṇyācyā praśnāvar āhe.}
\end{quote}

\begin{quote}
Today's editorial is about the water problem.
\end{quote}

Masculine.

\mr{अग्र}-, front, + \mr{लेख}: the leading article.

\begin{cousinweb}
\textbf{\mr{अग्रलेख}} (\emph{agralekh}) is a modern compound built from Sanskrit \emph{agra}, 'front, foremost', and \emph{lekha}, 'a writing' (Monier-Williams).
\end{cousinweb}

\begin{grammarlens}[title={What you've built}]
One more word for telling what happened.
\end{grammarlens}

\subsection*{Your turn}
\begin{itemize}
\raggedright
  \item \emph{Say it:} \emph{agralekh}
  \item \emph{Say it:} \emph{agralekh}, once more
  \item \emph{Say it:} \emph{ājcā agralekh pāṇyācyā praśnāvar āhe.} — Today's editorial is about the water problem.
  \item \emph{Recall:} say the Marathi for an editor, then the Marathi for an article, then say \emph{agralekh} again
\end{itemize}

\begin{tcolorbox}[breakable,colback=teal!4,colframe=teal!35!black,title={Before you move on}]
What does \mr{अग्रलेख} mean? (\textquotedblleft{}An editorial\textquotedblright{}.) Say it once more.
\end{tcolorbox}
\section[\texorpdfstring{\mr{प्रकाशित} \mr{करणे} (prakāśit karṇe)}{प्रकाशित करणे (prakāśit karṇe)}]{\mr{प्रकाशित} \mr{करणे} (prakāśit karṇe) — to publish}

\label{lesson:MR-C311-prakasit-karne}



\begin{quote}
Before the new one: say the Marathi for an article, then the Marathi for an editorial.
\end{quote}

\subsection*{You'll want to know: \mr{प्रकाशित} \mr{करणे}}
\textbf{\mr{प्रकाशित} \mr{करणे}} — \emph{prakāśit karṇe} — \textquotedblleft{}to publish\textquotedblright{}.

\begin{quote}
\mr{त्या} \mr{लेखकाने} \mr{नवं} \mr{पुस्तक} \mr{प्रकाशित} \mr{केलं}.
\end{quote}

\begin{quote}
\emph{tyā lekhkāne navaṁ pustak prakāśit kelaṁ.}
\end{quote}

\begin{quote}
That writer published a new book.
\end{quote}

Past, with -\mr{ने} (he …): \mr{त्याने} \mr{प्रकाशित} \mr{केलं}, \emph{tyāne prakāśit kelaṁ}.

What is published agrees; \mr{प्रकाशित} \mr{होणे}, to be published.

\begin{cousinweb}
\textbf{\mr{प्रकाशित} \mr{करणे}} (\emph{prakāśit karṇe}) joins \emph{prakāśit}, a tatsama taken from Sanskrit \emph{prakāśita}, 'brought to light, made public' (Monier-Williams), to \emph{karṇe}, 'to do'.
\end{cousinweb}

\begin{grammarlens}[title={What you've built}]
One more word for telling what happened.
\end{grammarlens}

\subsection*{Your turn}
\begin{itemize}
\raggedright
  \item \emph{Say it:} \emph{prakāśit karṇe}
  \item \emph{Say it:} \emph{prakāśit karṇe}, once more
  \item \emph{Say it:} \emph{tyā lekhkāne navaṁ pustak prakāśit kelaṁ.} — That writer published a new book.
  \item \emph{Recall:} say the Marathi for an article, then the Marathi for an editorial, then say \emph{prakāśit karṇe} again
\end{itemize}

\begin{tcolorbox}[breakable,colback=teal!4,colframe=teal!35!black,title={Before you move on}]
What does \mr{प्रकाशित} \mr{करणे} mean? (\textquotedblleft{}To publish\textquotedblright{}.) Say it once more.
\end{tcolorbox}
\section[(dialogue)]{Looking back: from Telling It Well to Who Makes the News}

\label{lesson:MR-R311-review-304-311}



\begin{quote}
Say the Marathi for to describe and an occasion.
\end{quote}

\subsection*{Your turn}
Say these aloud:

\begin{itemize}
\raggedright
  \item to describe, an occasion, to unfold, to make a stir, to spread
  \item last year, previous, before that, over time, to elapse
  \item step by step, many times, now and then, by itself, lately
  \item a commencement, an ending, at last, a period, a remembrance
  \item a generation, an ancestor, heritage, a lineage, to preserve
  \item to lose forever, to endure, death, to pass away, to arise
  \item a newspaper, a daily paper, a weekly, a headline, to print
  \item a reporter, an editor, an article, an editorial, to publish
\end{itemize}

\begin{tcolorbox}[breakable,colback=teal!4,colframe=teal!35!black,title={Before you move on}]
To describe? (\textbf{\mr{वर्णन} \mr{करणे}}, \emph{varṇan karṇe}.) An occasion? (\textbf{\mr{प्रसंग}}, \emph{prasaṅg}.) To unfold? (\textbf{\mr{उलगडणे}}, \emph{ulgaḍṇe}.) To make a stir? (\textbf{\mr{गाजणे}}, \emph{gājṇe}.) To spread? (\textbf{\mr{पसरणे}}, \emph{pasarṇe}.) Last year? (\textbf{\mr{गतवर्षी}}, \emph{gatvarṣī}.) Previous? (\textbf{\mr{मागील}}, \emph{māgīl}.) Before that? (\textbf{\mr{तत्पूर्वी}}, \emph{tatpūrvī}.) Over time? (\textbf{\mr{कालांतराने}}, \emph{kālāntarāne}.) To elapse? (\textbf{\mr{लोटणे}}, \emph{loṭṇe}.) Step by step? (\textbf{\mr{क्रमशः}}, \emph{kramaśaḥ}.) Many times? (\textbf{\mr{अनेकदा}}, \emph{anekdā}.) Now and then? (\textbf{\mr{अधूनमधून}}, \emph{adhūnmadhūn}.) By itself? (\textbf{\mr{आपोआप}}, \emph{āpoāp}.) Lately? (\textbf{\mr{नुकताच}}, \emph{nuktāc}.) A commencement? (\textbf{\mr{प्रारंभ}}, \emph{prārambh}.) An ending? (\textbf{\mr{अंत}}, \emph{ant}.) At last? (\textbf{\mr{अखेर}}, \emph{akher}.) A period? (\textbf{\mr{कालखंड}}, \emph{kālkhaṇḍ}.) A remembrance? (\textbf{\mr{स्मृती}}, \emph{smṛtī}.) A generation? (\textbf{\mr{पिढी}}, \emph{piḍhī}.) An ancestor? (\textbf{\mr{पूर्वज}}, \emph{pūrvaj}.) Heritage? (\textbf{\mr{वारसा}}, \emph{vārsā}.) A lineage? (\textbf{\mr{वंश}}, \emph{vaṁś}.) To preserve? (\textbf{\mr{जतन} \mr{करणे}}, \emph{jatan karṇe}.) To lose forever? (\textbf{\mr{गमावणे}}, \emph{gamāvṇe}.) To endure? (\textbf{\mr{सोसणे}}, \emph{sosṇe}.) Death? (\textbf{\mr{मृत्यू}}, \emph{mṛtyū}.) To pass away? (\textbf{\mr{मरण} \mr{पावणे}}, \emph{maraṇ pāvṇe}.) To arise? (\textbf{\mr{उद्भवणे}}, \emph{udbhavṇe}.) A newspaper? (\textbf{\mr{वृत्तपत्र}}, \emph{vṛttapatra}.) A daily paper? (\textbf{\mr{दैनिक}}, \emph{dainik}.) A weekly? (\textbf{\mr{साप्ताहिक}}, \emph{sāptāhik}.) A headline? (\textbf{\mr{मथळा}}, \emph{mathḷā}.) To print? (\textbf{\mr{छापणे}}, \emph{chāpṇe}.) A reporter? (\textbf{\mr{वार्ताहर}}, \emph{vārtāhar}.) An editor? (\textbf{\mr{संपादक}}, \emph{sampādak}.) An article? (\textbf{\mr{लेख}}, \emph{lekh}.) An editorial? (\textbf{\mr{अग्रलेख}}, \emph{agralekh}.) To publish? (\textbf{\mr{प्रकाशित} \mr{करणे}}, \emph{prakāśit karṇe}.)
\end{tcolorbox}
